\documentclass[11pt,a4paper]{ctexart}
\usepackage[utf8]{inputenc}
\usepackage{geometry}
\geometry{left=2.5cm,right=2.5cm,top=2.8cm,bottom=2.8cm}
\usepackage{amsmath,amssymb,amsfonts}
\usepackage{enumitem}
\usepackage{titlesec}
\usepackage{xcolor}
\usepackage{tcolorbox}
\usepackage{hyperref}

% 引入画图和绝对位置放置宏包用于生成水印
\usepackage{graphicx}
\usepackage{eso-pic}
\usepackage{tikz} % 引入 TikZ 宏包用于绘图

% 解决带圈数字在部分编译器下报错的问题
\xeCJKDeclareCharClass{CJK}{"2460->"2469}

% --- 颜色定义 ---
\definecolor{mainblue}{RGB}{26, 82, 118}
\definecolor{secblue}{RGB}{41, 128, 185}
\definecolor{bglight}{RGB}{244, 246, 247}
\definecolor{textdark}{RGB}{44, 62, 80}

% --- 页面美化与字体设置 ---
\hypersetup{colorlinks=true, linkcolor=mainblue, urlcolor=secblue}
\setlist[itemize]{leftmargin=2em, itemsep=0.3em}
\setlist[enumerate]{leftmargin=2em, itemsep=0.3em}

% --- 标题样式自定义 ---
\titleformat{\section}{\large\bfseries\color{mainblue}}{\thesection}{1em}{}[{\color{mainblue}\titlerule[1pt]}]
\titleformat{\subsection}{\normalsize\bfseries\color{secblue}}{\thesubsection}{1em}{}
\titleformat{\subsubsection}{\small\bfseries\color{textdark}}{\thesubsubsection}{1em}{}

% --- 水印设置 (每页三个，斜45度，居中分布) ---
\newcommand{\WMText}{贺昱琪学习资料}
\newcommand{\WMScale}{2.28}
\newcommand{\WMAngle}{45}
\colorlet{WMColor}{black!8}

\AddToShipoutPictureBG{%
  \AtPageLowerLeft{%
    \put(\LenToUnit{0.66\paperwidth},\LenToUnit{0.70\paperheight}){%
      \rotatebox{\WMAngle}{\scalebox{\WMScale}{\textcolor{WMColor}{\WMText}}}%
    }%
    \put(\LenToUnit{0.33\paperwidth},\LenToUnit{0.50\paperheight}){%
      \rotatebox{\WMAngle}{\scalebox{\WMScale}{\textcolor{WMColor}{\WMText}}}%
    }%
    \put(\LenToUnit{0.66\paperwidth},\LenToUnit{0.30\paperheight}){%
      \rotatebox{\WMAngle}{\scalebox{\WMScale}{\textcolor{WMColor}{\WMText}}}%
    }%
  }%
}

\begin{document}

\begin{center}
    \LARGE \textbf{同类型练习题（贺昱琪）}
\end{center}

\vspace{0.5cm}

\noindent \textbf{第一类：统计与平均数综合问题（类似原第8题）}

\begin{enumerate}
    \item 某班全体学生进行了一次投篮测试，每人投 10 个。得 0 分的有 3 人，1 分的有 2 人，2 分的有 5 人；得 8 分的有 4 人，9 分的有 2 人，10 分的有 1 人；其余分数人数未知。已知该班学生中，至少得 3 分的人平均得 7 分，得分不到 8 分的人平均得分为 4 分。那么该班共有\underline{\hspace{2cm}}人。\vspace{1.2cm}
    
    \item 一次 10 分制的答题比赛中，得 0 分的有 5 人，1 分的有 4 人，2 分的有 1 人；得 9 分的有 3 人，10 分的有 2 人。已知至少得 3 分的人平均得 8 分，不到 9 分的人平均得 5 分。问参加比赛的共有\underline{\hspace{2cm}}人。\vspace{1.2cm}
    
    \item 某小组参加满分 10 分的操作考核。得 0 分 2 人，1 分 3 人，2 分 3 人；得 7 分 2 人，8 分 2 人，9 分 1 人，10 分 1 人。已知至少得 3 分的人平均得 6 分，不到 7 分的人平均得 4 分。求该小组的总人数是\underline{\hspace{2cm}}人。\vspace{1.2cm}
    
    \item 某次满分为 10 分的随堂测验中，得 0 分 1 人，1 分 1 人，2 分 1 人，3 分 2 人；得 8 分 3 人，9 分 1 人，10 分 1 人。已知至少得 4 分的人平均得 7 分，不到 8 分的人平均得 4 分。问该班共有\underline{\hspace{2cm}}人。\vspace{1.2cm}
    
    \item 10 分制的体能测试中，得 0 分 4 人，1 分 3 人，2 分 3 人；得 9 分 3 人，10 分 2 人。已知至少得 3 分的人平均得 8 分，不到 9 分的人平均得 5 分。请问参加测试的共有\underline{\hspace{2cm}}人。\vspace{1.2cm}
\end{enumerate}

\vspace{0.8cm}

\noindent \textbf{第二类：工程与工期问题（类似原第14题）}

\begin{enumerate}
    \item 单独完成一项任务，甲按规定时间可提前 4 天完成，乙则要超过规定时间 6 天才能完成。如果甲、乙二人合作 4 天后，剩下的由乙单独做，那么刚好在规定时间完成。问：甲、乙二人合作需多少天完成该任务？
    \vspace{3.5cm}
    
    \item 单独完成一项任务，甲按规定时间可提前 2 天完成，乙则要超过规定时间 4 天才能完成。如果甲、乙二人合作 2 天后，剩下的由乙单独做，那么刚好在规定时间完成。问：甲、乙二人合作需多少天完成？
    \vspace{3.5cm}
    
    \item 单独完成一项任务，甲按规定时间可提前 3 天完成，乙则要超过规定时间 5 天才能完成。如果甲、乙二人合作 3 天后，剩下的由乙单独做，那么刚好在规定时间完成。问：甲、乙二人合作需多少天完成？
    \vspace{3.5cm}
    
    \item 单独完成一项任务，甲按规定时间可提前 10 天完成，乙则要超过规定时间 20 天才能完成。如果甲、乙二人合作 5 天后，剩下的由乙单独做，那么刚好在规定时间完成。问：甲、乙二人合作需多少天完成？
    \vspace{3.5cm}
    
    \item 单独完成一项任务，甲按规定时间可提前 10 天完成，乙则要超过规定时间 30 天才能完成。如果甲、乙二人合作 10 天后，剩下的由乙单独做，那么刚好在规定时间完成。问：甲、乙二人合作需多少天完成？
    \vspace{3.5cm}
\end{enumerate}

\vspace{0.8cm}

\noindent \textbf{第三类：成本、售价与利润率问题（类似原第16题）}

\begin{enumerate}
    \item 某工厂生产一款产品的成本随着物价上涨比原来增加了 20\%，为了保证收益，厂家把售价提高了 50\%（即 0.5 倍），最终核算发现利润率比原来增加了 30\%。求该产品原来的利润率。
    \vspace{3.5cm}
    
    \item 某开发商开发一套房子的成本比原来增加了 25\%，为了赚钱，开发商把售价提高了 75\%（即 0.75 倍），结果利润率比原来增加了 60\%。求开发商原来的利润率。
    \vspace{3.5cm}
    
    \item 某电子产品的进货成本比原来增加了 10\%，零售商顺势把售价提高了 32\%，最终该产品的利润率比原来增加了 28\%。求该产品原来的利润率。
    \vspace{3.5cm}
    
    \item 由于原材料紧缺，某商品的成本增加了 50\%，商家将售价提高了 80\%，结算时发现利润率比原来增加了 25\%。求该商品原来的利润率。
    \vspace{3.5cm}
    
    \item 某餐饮店的食材成本比原来增加了 5\%，老板将对应菜品的价格提高了 26\%，计算后发现该菜品的利润率比原来增加了 32\%。求该菜品原来的利润率。
    \vspace{3.5cm}
\end{enumerate}

\end{document}